\documentclass[../veristore_security_model.tex]{subfiles}

\begin{document}
\section{Trust Levels}

% COMMENT: Define the hierarchy of trust levels in the system. For each trust level, describe
% what entities operate at that level and what capabilities/privileges they have. In distributed
% systems, this often includes: system administrator, authenticated client, storage server,
% untrusted network attacker. Note that some components may operate at different trust levels
% depending on deployment scenario.

\subsection{Administrator}

% COMMENT: Describe the system administrator trust level. This includes anyone with physical or
% root access to the servers. Administrators have complete control over the system and can read,
% modify, or delete any data. Security model assumes administrators are trusted and benign.
System administrators have full access to the servers hosting the veri-store system and the stored data. They are responsible for maintaining the infrastructure, applying updates, and ensuring the system is running smoothly. The security model assumes that administrators are trusted and will not intentionally compromise data integrity or confidentiality. However, they have the capability to access stored fragments and fingerprints, so it is important to consider this trust level when designing mitigations for information disclosure threats.
\vspace{1em}

\noindent Example System Administrators:
\begin{itemize}
    \item Developers (e.g., Albert and Mac) during development and testing
    \item IT staff responsible for server/data maintenance in a production deployment
    \item Cloud provider administrators if deployed on cloud infrastructure (e.g., AWS, Azure)
    \item System operators with root access to physical or virtual servers
    \item \textit{Additional examples to be added}
\end{itemize}

\subsection{Authenticated Client}
The term ``authenticated client'' refers both to the end users and the application they use to access the Veri-Store system.

An authenticated client has verified their identity to the system and is authorized to store and retrieve their own data. They are assumed to be legitimate, but may make mistakes or run buggy software which can violate system security. They should not have access to other users' data or have the ability to violate system integrity.

\subsection{Storage Server}

% COMMENT: Describe the trust level of storage servers. In the semi-trusted model, servers might
% fail, return corrupted data, or attempt to forge fingerprints, but do not actively collude
% (beyond the f-threshold). Specify whether servers are assumed to follow protocols correctly or
% might exhibit Byzantine behavior.

\subsection{Untrusted Network}
We assume a TLS connection between clients and servers, although we do not rely on it for data integrity. Adversaries will be able to observe network traffic, and it is assumed that they will, in theory, be able to inject network traffic. 

However, the homomorphic fingerprinting scheme allows for any inconsistencies between data fragments to be caught before reconstruction.  Any tampering will be detected because the tampered fragment's fingerprint won't match the expected encoding of the other fragments' fingerprints. In order to not be caught, an adversary would need to provide a consistent set of tampered fragments which all pass verification. This is computationally infeasible given the collision resistance of the fingerprinting scheme.

Unfortunately, given the material constraints of this project (i.e., ``servers'' existing on a single machine as different \texttt{localhost} ports), we are unable to fully test this.

\newpage
\end{document}